% =====================================================================
%  BODY — Band VII · Der Markt der Begehrten
% =====================================================================

\section{Die Umkehrung der Leitfrage}

Die bisherigen sechs Bände trugen eine einzige Formel durch drei Rechtskreise
und eine biologische Koda: \emph{Rechtsstellung folgt Funktionsstellung.} Wer
eine gesellschaftlich unverzichtbare Funktion innehatte, erhielt Rechte, die
diese Funktion absicherten; wer sie nicht innehatte, wurde rechtlich gemindert.
Band VI wandte die Formel erstmals auf den Mann an und zeigte, dass auch seine
scheinbaren Privilegien Funktionsprämien waren --- Prämien für Wehrdienst,
Haftung und Ernährung, entziehbar bei Funktionsverlust.

Dieser Band verschiebt den Schauplatz von der Rechtsordnung in ihr intimstes
Vorzimmer: die Partnerwahl. Die These lautet, dass der moderne Partnermarkt
kein Gegenort zur Funktionslogik ist, sondern ihr letzter, kaum regulierter
Vollzugsraum. Wo das Recht die Funktion nicht mehr erzwingt --- der Staat
schreibt keinem Mann mehr vor, Ernährer zu sein ---, erzwingt sie womöglich die
Selektion. Und sie erzwingt sie nicht durch Gesetz, sondern durch etwas
Wirkmächtigeres: durch Begehren.

\begin{merksatz}
Wenn das Recht die Funktion freigibt und das Begehren sie festhält, dann ist
die Partnerwahl der Ort, an dem das Patriarchat sich selbst reproduziert ---
ohne Gesetzgeber, ohne Zwang, scheinbar aus freien Stücken.
\end{merksatz}

Das ist eine unbequeme These, weil sie die Rollen verschiebt. Sie fragt nicht
mehr nur, wie Männer Frauen an Funktionen banden, sondern ob die Selektion der
Begehrten ihrerseits Männer an Funktionen bindet --- und ob die Freiheit der
Wahl, die die westliche Kultur als ihre Errungenschaft feiert, auf einer Seite
des Marktes zur Freiheit wird, auf der anderen zur Prüfung. Der Band verfolgt
diese Frage in fünf Schritten: erst die empirische Vermessung des Marktes
(§\,2--3), dann das theoretische Herzstück der Ambivalenz (§\,4), dann die
Ökonomie des Mangels (§\,5), dann der Offline-Rest --- die unvermittelte
Partnerwahl und ihre eigenen Regeln (§\,6) ---, schließlich die Ausstiegsfrage
(§\,7) und die Synthese (§\,8).

Eine Vorbemerkung zur Redlichkeit: Dieser Gegenstand ist vermint. Er wird
rechts von einer Manosphere bewirtschaftet, die aus Selektionsdaten
Geschlechterhass destilliert, und links von einer Reflexabwehr, die jede
Thematisierung männlichen Leidens für einen Angriff auf feministische
Errungenschaften hält. Beide Lager teilen einen Fehler: Sie lesen Asymmetrie
als Schuld. Dieser Band liest sie als Struktur. Er wird die Daten so hart
nehmen, wie sie sind, und die Schuldzuweisung so konsequent verweigern, wie es
die Struktur verlangt.

\section{Die Vermessung des Marktes I: Bewertung und Begehren}

\subsection{Die Bewertungsasymmetrie}

Der meistzitierte Befund stammt aus der internen Datenauswertung der
Dating-Plattform OkCupid, veröffentlicht 2014 vom Mitgründer Christian Rudder
in \textit{Dataclysm}. \belegt\ Auf die Frage, wie Nutzer die Attraktivität des
jeweils anderen Geschlechts bewerten, ergaben sich zwei fundamental
verschiedene Kurven. Männer bewerteten die Attraktivität von Frauen annähernd
normalverteilt --- die Masse landete in der Mitte, wenige sehr hoch, wenige sehr
niedrig, das erwartbare Bild einer Glockenkurve. Frauen hingegen bewerteten die
Männer stark linksschief: Der Gipfel der Verteilung lag deutlich unterhalb der
Mitte. Rund vier Fünftel der Männer wurden als unterdurchschnittlich attraktiv
eingestuft.

\begin{datenbox}[BEFUND · OKCUPID / RUDDER 2014]
Frauen bewerteten ca.~80\,\% der Männer als \emph{unterdurchschnittlich}
attraktiv. Männer verteilten ihre Bewertungen von Frauen annähernd
symmetrisch um den Mittelwert. Rudders eigene Zuspitzung: Übertrüge man diese
Wahrnehmung auf Intelligenz, entspräche sie einer Welt, in der die Mehrheit der
Frauen die Mehrheit der Männer für kognitiv beeinträchtigt hielte.
\end{datenbox}

Zur methodischen Einordnung: Es handelt sich um Plattform-Verhaltensdaten einer
Selbstauswahl-Population (US-amerikanische OkCupid-Nutzer um 2009--2013), nicht
um eine repräsentative Bevölkerungsstichprobe. \umstritten\ Die Zahl ist als
Wahrnehmungsartefakt robust und wurde von einem deutschen Anbieter mit eigenen
Daten annähernd repliziert (~81\,\%), aber sie misst \emph{geäußerte
Attraktivitätsurteile im Online-Kontext}, nicht Partnerpräferenzen im Leben.
Wer sie als Aussage über \enquote{die Frauen} liest, überdehnt sie. Wer sie als
Strukturmerkmal des Online-Marktes liest, trifft sie.

\subsection{Das aspirationale Prinzip --- und warum es nicht weiblich ist}

Der häufigste Fehlschluss aus der Bewertungsasymmetrie lautet: Frauen
\enquote{daten nach oben}, Männer \enquote{nach unten}. Die belastbarste Studie
zur Kontaktdynamik widerspricht der geschlechtlichen Zuspitzung. Elizabeth
Bruch und Mark Newman analysierten 2018 in \textit{Science Advances} das
vollständige Nachrichtennetzwerk einer großen Plattform in vier US-Metropolen
und rekonstruierten daraus eine hierarchische Begehrtheitsordnung. \belegt

\begin{datenbox}[BEFUND · BRUCH \& NEWMAN 2018, SCIENCE ADVANCES]
\textbf{Beide} Geschlechter schreiben bevorzugt Personen an, die im Mittel rund
25\,\% \emph{begehrter} sind als sie selbst. Das aspirationale Anschreiben ist
also \emph{kein} weibliches Muster, sondern ein allgemein menschliches.
Zusätzlich: Die Begehrtheit von Frauen erreicht ihren Gipfel bei etwa 18--22
Jahren und sinkt danach kontinuierlich; die Begehrtheit von Männern steigt
langsamer, gipfelt erst um die 50 und ist stärker an Bildung und Einkommen
gekoppelt.
\end{datenbox}

Der entscheidende Punkt für diesen Band liegt in der Kombination beider
Befunde. Nicht das \emph{Wollen} unterscheidet die Geschlechter --- beide wollen
nach oben ---, sondern die \emph{Marktmacht}, dieses Wollen einzulösen. Weil die
weibliche Aufmerksamkeit sich (siehe §\,2.1) auf ein schmales Männersegment
konzentriert, kann die durchschnittliche Frau ihre 25\,\%-Aspiration real
verwirklichen: Das begehrte Segment schreibt zurück. Der durchschnittliche Mann
kann es nicht: Sein aspirationales Anschreiben verhallt, weil die von ihm
angeschriebenen Frauen bereits vom begehrten Männersegment umworben werden.
Dieselbe Strategie führt auf der einen Seite zum Erfolg, auf der anderen ins
Leere --- nicht wegen unterschiedlicher Absichten, sondern wegen
unterschiedlicher struktureller Position.

\begin{merksatz}
Aspirationales Daten ist geschlechtsneutral. Was sich unterscheidet, ist die
Wahrscheinlichkeit, dass die Aspiration erwidert wird. Die Frau daatet nicht
\enquote{gieriger} nach oben --- sie hat nur die Marktmacht, oben auch gehört zu
werden.
\end{merksatz}

Das ist analytisch die sauberere Formulierung und zugleich die, die der
Manosphere den Boden entzieht: Es gibt kein moralisches Defizit weiblicher
Wahl, es gibt eine strukturelle Asymmetrie der Marktposition. Genau diese
Asymmetrie ist aber real und für den Durchschnittsmann folgenreich --- und sie
kleinzureden wäre der spiegelbildliche Fehler.

\section{Die Vermessung des Marktes II: Match-Ökonomie und Ungleichheit}

\subsection{Die Match-Raten}

Die früheste systematische Erhebung der Tinder-Dynamik legten Gareth Tyson und
Kollegen 2016 vor (\textit{A First Look at User Activity on Tinder}). \belegt\
Mit kontrollierten Profilen maßen sie das Wisch- und Match-Verhalten.

\begin{datenbox}[BEFUND · TYSON ET AL. 2016]
Männliche Profile erhielten in der Erhebung eine Match-Rate von rund
\textbf{0,6\,\%}; weibliche Profile lagen um \textbf{10\,\%} --- etwa ein
Faktor 15. Männer wischen dabei wenig selektiv (großzügiges Rechtswischen),
Frauen hochselektiv. Die männliche Strategie des Massenwischens ist eine
rationale Antwort auf die niedrige Trefferquote --- und verschärft zugleich die
Übersättigung der begehrten weiblichen Profile.
\end{datenbox}

Nachfolgende Datensätze bestätigen die Größenordnung, nicht die exakte Zahl:
Männer wischen bei grob 40--60\,\% der Profile nach rechts, Frauen bei rund
5--15\,\%, und trotzdem --- oder gerade deswegen --- erhalten Frauen ein
Vielfaches an Matches. \plausibel\ Die Zahlen schwanken je nach Plattform,
Region und Erhebungsjahr; als \emph{Ordnung} ist die Asymmetrie durchgängig.

\subsection{Der Gini-Koeffizient des Begehrens}

Um die Konzentration der Aufmerksamkeit zu quantifizieren, haben Beobachter den
Gini-Koeffizienten herangezogen --- jenes Ungleichheitsmaß, das sonst
Einkommensverteilungen beschreibt (0 = alle gleich, 1 = einer hat alles).

\begin{datenbox}[BEFUND · LIKE-KONZENTRATION (mehrere Quellen)]
Eine vielzitierte Tinder-Analyse schätzte den Gini der männlichen
\enquote{Attraktivität} (gemessen an erhaltenen Likes) auf rund \textbf{0,58}
--- ungleicher als die Einkommensverteilung der meisten Volkswirtschaften. Eine
Auswertung von Hinge-Daten (Aviv Goldgeier) ergab ~\textbf{0,54} für Männer
gegenüber ~\textbf{0,38} für Frauen: Die von Männern erhaltene Aufmerksamkeit
ist also weit ungleicher verteilt als die von Frauen erhaltene.
\end{datenbox}

\begin{gegen}
Diese Gini-Werte stammen aus Hobby- und Firmenanalysen, nicht aus
peer-reviewter Forschung. Die bekannteste Tinder-Zahl (0,58) beruht auf einer
kleinen Zahl konstruierter Profile und einer Hochrechnung, nicht auf einer
Vollerhebung. Sie taugt als \emph{Illustration} der Konzentration, nicht als
präziser Messwert. Wer mit ihr wie mit einer amtlichen Statistik argumentiert,
überschreitet, was die Datenbasis trägt. Als Größenordnung fügt sie sich
allerdings konsistent in die peer-reviewten Befunde von Tyson sowie Bruch \&
Newman.
\end{gegen}

\subsection{Was die Konzentration bedeutet}

Fasst man die drei Befunde zusammen --- linksschiefe Bewertung (§\,2.1),
aspirationale Erwiderung nur für Begehrte (§\,2.2), hochkonzentrierte Matches
(§\,3.1--3.2) ---, ergibt sich das Bild eines Marktes mit extremer
Aufmerksamkeitskonzentration auf der weiblichen \emph{und} der männlichen
Spitzenseite, bei struktureller Unsichtbarkeit der männlichen Mitte und
Unterseite. \plausibel\ Der durchschnittliche Mann ist auf diesem Markt nicht
\emph{abgelehnt} --- er ist \emph{unsichtbar}, was subjektiv schwerer wiegt, weil
Ablehnung wenigstens Wahrnehmung voraussetzt.

\subsection{Die Grenzen der Selbstauskunft}

Eine methodische Warnung gehört in die Vermessung selbst: Profilangaben sind
keine Messwerte, sondern Selbstdarstellung. Die klassische
Kreuzvalidierungsstudie von Toma, Hancock und Ellison (2008) verglich
Profilangaben mit nachgemessenen Werten: \textbf{81\,\%} der Online-Dater
wichen bei Größe, Gewicht oder Alter messbar von der Wahrheit ab. \belegt\
Männer übertrieben systematisch die Größe, Frauen untertrieben das Gewicht;
die Fotos waren das am stärksten geschönte Element. Die Abweichungen waren
klein --- kalkuliert unterhalb der Nachweisbarkeitsschwelle des ersten Dates.

Folgenreicher als die Zentimeter ist die \emph{Absichtsangabe}.
\enquote{Suche feste Beziehung} ist billig zu behaupten und teuer zu prüfen ---
in der Sprache der Spieltheorie: cheap talk. Die Täuschung läuft in beide
Richtungen. Der Mann, der Bindungswunsch angibt und Unverbindliches meint, ist
das eine Muster; experimentell ist belegt, dass die Bereitschaft, über eigene
Absichten, Einkommen und Persönlichkeit zu täuschen, mit der Attraktivität des
Gegenübers steigt --- \emph{ohne Geschlechtsunterschied} in der
Täuschungsbereitschaft (Rowatt et al.). \plausibel\ Das Gegenstück: die Frau,
die unverbindliches Interesse \emph{untertreibt}, weil die soziale
Sanktionierung weiblicher Promiskuität die ehrliche Angabe verteuert.
\plausibel\ Beide Verzerrungen wirken in dieselbe Richtung: Die geäußerten
Absichten überzeichnen die Bindungsorientierung des Marktes.

\begin{merksatz}
Profilangaben messen, was Menschen signalisieren wollen --- nicht, was sie
suchen. Deshalb ruht die Beweisführung dieses Bandes auf Verhaltensdaten:
Matches, Nachrichten, Trennungen. Die Absichtsstatistiken (§\,6) sind mit
diesem Vorbehalt zu lesen.
\end{merksatz}

Damit ist die empirische Grundlage gelegt. Die eigentliche Frage dieses Bandes
ist keine der Zahlen, sondern der Deutung: Was folgt daraus für die Schuld- und
Systemfrage? Dem widmet sich das Herzstück.

\section{Das Herzstück: Reproduzieren die Begehrten das Patriarchat?}

Hier liegt die eigentliche Provokation des Bandes, und sie muss präzise gestellt
werden, um nicht in Schuldzuweisung zu kippen. Die Frage lautet nicht:
\enquote{Sind Frauen schuld am Patriarchat?} Sie lautet: \emph{Wenn} weibliche
Partnerwahl im Aggregat Status, Einkommen und Ernährerfähigkeit belohnt --- und
die Empirie in §\,2--3 legt nahe, dass sie es im Mittel tut ---, stabilisiert
diese Wahl dann nicht genau jene Funktionslogik, die den Mann an die
Ernährerrolle kettet und die Frau an ihre Kehrseite?

\subsection{Bourdieu: Die Beherrschten vollstrecken die Herrschaft}

Die tragfähigste Antwort liefert Pierre Bourdieus \textit{Die männliche
Herrschaft} (1998). \plausibel\ Sein Begriff der \emph{symbolischen Gewalt}
beschreibt eine Herrschaft, die nicht von außen erzwungen wird, sondern in die
Wahrnehmungs- und Bewertungsschemata der Beherrschten selbst eingelassen ist.
Die Beherrschten nehmen die Welt mit Kategorien wahr, die das
Herrschaftsverhältnis hervorgebracht hat --- und reproduzieren es dadurch in
scheinbar freien Akten. Wer Status begehrt, hat die Statusordnung bereits
verinnerlicht. Die Wahl ist frei \emph{und} vorgeformt zugleich; das ist kein
Widerspruch, sondern die Funktionsweise symbolischer Gewalt.

Auf den Partnermarkt übertragen: Die Präferenz für den ressourcenstarken Mann
ist keine \enquote{Berechnung} und kein \enquote{Verrat}, sondern eine tief
sozialisierte Disposition, die die Statusordnung fortschreibt, gerade weil sie
sich als persönlicher Geschmack erlebt. \enquote{Er wirkt souverän},
\enquote{er hat sein Leben im Griff}, \enquote{ich will jemanden, zu dem ich
aufschauen kann} --- das sind die Alltagsformeln, in denen sich die Statusordnung
als Gefühl tarnt.

\begin{merksatz}
Die Begehrten sind nicht die Architektinnen der Ordnung. Sie sind ihre
Vollstreckerinnen --- so wie die Begehrenden ihre Antragsteller sind. Beide
handeln frei innerhalb einer Ordnung, die keiner von beiden entworfen hat.
\end{merksatz}

\subsection{Kandiyoti: Der patriarchale Handel}

Deniz Kandiyotis Konzept des \emph{patriarchal bargain} (1988) schärft das
Bild. \plausibel\ Frauen agieren, so Kandiyoti, rational innerhalb der ihnen
gesetzten Strukturen: Sie maximieren Sicherheit und Lebenschancen unter
patriarchalen Bedingungen --- und stabilisieren dadurch die Struktur, die diese
Bedingungen setzt. Wer in einer Welt lebt, in der männliches Einkommen
verlässlicher, höher und karrierestabiler ist (Band VI dokumentierte, warum),
für den ist die Präferenz für einen ressourcenstarken Partner keine
Ideologie, sondern eine vernünftige Risikoabsicherung. Die
\enquote{Wahl} ist eine Anpassung an reale Anreize.

Das erklärt auch, warum die Präferenz sich nicht einfach \enquote{wegerziehen}
lässt: Solange die zugrunde liegende ökonomische Asymmetrie besteht, bleibt der
Handel rational. Man kann Menschen schwer vorwerfen, dass sie auf reale Anreize
reagieren.

\subsection{Der Test: Verschwindet die Präferenz mit der Gleichstellung?}

Damit ist die entscheidende empirische Frage gestellt: Ist die weibliche
Statuspräferenz ein \emph{Anpassungsprodukt} (dann müsste sie mit
zunehmender ökonomischer Gleichstellung verschwinden) oder eine \emph{stabile
Disposition} (dann bliebe sie auch bei Gleichstellung)? Hier stehen sich zwei
Forschungsprogramme gegenüber.

\textbf{Sozialkonstruktivistische Position (Eagly \& Wood, Social Role
Theory).} \umstritten\ Partnerpräferenzen folgen der geschlechtlichen
Arbeitsteilung. Wo Frauen ökonomisch abhängig sind, präferieren sie Ressourcen;
wo die Abhängigkeit sinkt, sollte auch die Präferenz sinken. Eagly und Wood
reanalysierten Buss' berühmte 37-Kulturen-Daten und fanden: Je gleichstellter
ein Land, desto schwächer die geschlechtstypischen Präferenzunterschiede. Das
stützt die Anpassungsthese.

\textbf{Evolutionspsychologische Position (Buss, Schmitt).} \umstritten\ Die
Statuspräferenz sei ein evolutionär geformtes, kulturübergreifend stabiles
Muster. Buss' Daten über 37 Kulturen zeigen die Ressourcenpräferenz von Frauen
in \emph{jeder} untersuchten Kultur --- variabel in der Stärke, konstant in der
Richtung.

\begin{gegen}
Die Datenlage stützt keine der beiden Positionen eindeutig, und genau das ist
redlich zu benennen. Manche Studien finden eine Abschwächung der
Statuspräferenz in egalitäreren Ländern; andere finden das
\enquote{Gender-Equality-Paradox}: In den gleichstellten skandinavischen
L�ndern werden manche geschlechtstypischen Präferenzunterschiede nicht kleiner,
sondern \emph{größer}. Speed-Dating-Studien (Eastwick \& Finkel 2008) zeigen
zudem eine Lücke zwischen \emph{genannten} und \emph{gelebten} Präferenzen:
Was Menschen als Kriterium angeben, sagt ihr tatsächliches Wahlverhalten oft
schlecht vorher. Wer hier eine der beiden Seiten zur gesicherten Wahrheit
erklärt, verkauft eine offene Forschungskontroverse als entschieden.
\end{gegen}

Für diesen Band ist die Nicht-Entscheidung produktiv. Denn beide Positionen
führen zur selben strukturellen Konsequenz: Ob die Statuspräferenz nun
anerzogen oder angelegt ist --- \emph{solange sie wirkt}, bindet sie den Mann an
die Funktion. Die Ursachenfrage ändert nichts an der Wirkung. Und die Wirkung
ist der Gegenstand dieses Bandes.

\subsection{Sind wir nicht alle Gefangene des Systems?}

Damit ist die Täter-Opfer-Dichotomie unterlaufen. Wenn die Begehrten die
Ordnung vollstrecken, ohne sie entworfen zu haben, und die Begehrenden sich ihr
unterwerfen, ohne sie gewählt zu haben --- wer ist dann der Nutznießer?

Bourdieu gibt die radikalste Antwort: Auch der Herrschende ist von seiner
Herrschaft beherrscht. Der Mann, der Status verkörpern muss, um begehrt zu
werden, ist nicht frei --- er ist zur permanenten Leistung verurteilt (die
\emph{Precarious Manhood} aus Band VI). bell hooks formulierte es in
\textit{The Will to Change} (2004) für die Gegenseite: Das Patriarchat
verstümmelt auch Männer, indem es ihnen emotionale Ganzheit im Tausch gegen
Macht abkauft. \plausibel

\begin{merksatz}
Es gibt auf diesem Markt keine reinen Gewinner. Es gibt eine schmale begehrte
Spitze beider Geschlechter, eine große funktional gebundene Mitte und eine
unsichtbare Unterseite. Das ist keine Ordnung von Tätern und Opfern. Es ist
eine Ordnung, in der fast alle zugleich vollstrecken und erleiden.
\end{merksatz}

\begin{gegen}
Die Systemperspektive darf nicht zur Verharmlosung realer Asymmetrien werden.
Dass \enquote{alle leiden}, ist wahr und zugleich gefährlich, wenn es die
fortbestehenden strukturellen Vorteile von Männern einebnet: den Gender Pay
Gap, die ungleiche Verteilung der Sorgearbeit, die überwältigende männliche
Mehrheit bei schwerer Partnergewalt, die gläserne Decke. Die
\enquote{patriarchale Dividende} (Connell) ist real; sie wird nur nicht von
allen Männern eingelöst. Wer aus \enquote{alle sind Gefangene} ein
\enquote{also sind Männer die eigentlichen Opfer} macht, betreibt genau die
Verkehrung, vor der dieser Band warnt. Symmetrie in der Analyse bedeutet nicht
Gleichheit der Lasten.
\end{gegen}

\section{Die Ökonomie des Mangels: Hypergamie und Monetarisierung}

\subsection{Hypergamie jenseits der App}

Der Verdacht, die App verzerre nur, was offline nicht existiere, hält der
Empirie nicht stand. Die Statuskopplung zeigt sich auch in den härtesten
Offline-Daten --- den Trennungsstatistiken.

\begin{datenbox}[BEFUND · KILLEWALD 2016, AM. SOCIOLOGICAL REVIEW]
In einer Analyse US-amerikanischer Ehen erhöhte die \emph{Erwerbslosigkeit des
Mannes} das Trennungsrisiko deutlich --- die Erwerbslosigkeit der Frau tat dies
nicht. Nicht die relative Einkommensverteilung war ausschlaggebend, sondern die
Verletzung der männlichen Ernährernorm: Ein Mann ohne Erwerbsarbeit destabilisiert
die Ehe messbar, eine Frau ohne Erwerbsarbeit nicht.
\end{datenbox}

\begin{datenbox}[BEFUND · BERTRAND, KAMENICA \& PAN 2015, QJE]
Sobald die Ehefrau mehr verdiente als der Ehemann, sank die Heiratsneigung, stieg
die Unzufriedenheit, wuchs die Scheidungswahrscheinlichkeit --- und paradoxerweise
übernahm die besserverdienende Frau \emph{mehr} Hausarbeit, offenbar um die
Normverletzung symbolisch zu kompensieren. Die Verteilung um die
50/50-Einkommensgrenze zeigte einen sichtbaren Knick: Paare vermeiden die
Konstellation \enquote{Frau verdient mehr}.
\end{datenbox}

Diese Befunde sind für die Kernthese entscheidend. \belegt\ Sie zeigen, dass die
Ernährernorm nicht ein Artefakt der Wischlogik ist, sondern eine tief verankerte
Strukturerwartung, die noch in modernen, formal gleichgestellten Ehen
wirksam ist und Verhalten steuert --- bis in die Aufteilung der Hausarbeit
hinein. Der Mann wird nicht durch die App zum Ernährer gemacht; die App macht
nur sichtbar und messbar, was die Ehe verdeckt fortführt.

\begin{gegen}
Zur Präzision: Die Bildungshypergamie hat sich umgekehrt. In den meisten
westlichen Ländern heiraten Frauen heute häufiger Männer mit \emph{niedrigerer}
formaler Bildung als sie selbst --- schlicht, weil Frauen die
Bildungsabschlüsse dominieren (Esteve et al.). Die \emph{Einkommens}- und
\emph{Status}hypergamie ist damit aber nicht verschwunden; sie hat sich nur von
der Bildung entkoppelt. Frauen heiraten seltener \enquote{nach oben} im
Diplom, aber weiterhin überwiegend \enquote{nicht nach unten} im Einkommen und
im wahrgenommenen Status. Die Norm ist nicht tot, sie hat ihre Währung
gewechselt.
\end{gegen}

\subsection{Die Zahlsklaven-These}

Wenn ein Markt eine strukturell unterlegene Seite hat, entsteht ein
Geschäftsmodell, das ihren Mangel bewirtschaftet. Genau das leisten die
Monetarisierungsmechaniken der Dating-Apps.

\begin{datenbox}[BEFUND · MARKTSTRUKTUR]
Rund drei Viertel der Tinder-Nutzerschaft sind männlich. \plausibel\ Die
kostenpflichtigen Sichtbarkeits- und Reichweitenprodukte (Boost, Super-Like,
Premium-Stufen) sind strukturell Produkte für die Marktseite mit
Aufmerksamkeitsdefizit --- also überwiegend für Männer. Die Match Group
erwirtschaftet ihren Umsatz damit überproportional aus männlichen zahlenden
Nutzern.
\end{datenbox}

Die eigentliche Pointe ist eine perverse: Auswertungen deuten darauf hin, dass
gerade die Männer, die am aggressivsten in Sichtbarkeit investieren (viele
Super-Likes kaufen), \emph{unterdurchschnittliche} Erfolgsraten haben.
\plausibel\ Das Geschäftsmodell verkauft dem funktional Unterlegenen die
Illusion eines Auswegs, ohne den Ausweg zu liefern --- denn ein tatsächlich
gelöstes Matching-Problem wäre ein verlorener Kunde.

\begin{merksatz}
Die Plattform hat kein ökonomisches Interesse an deinem Erfolg. Sie hat ein
Interesse an deinem Mangel --- weil der Mangel das Abonnement verlängert. Der
optimale Kunde ist der ewig Suchende.
\end{merksatz}

Damit schließt sich der Kreis zur Funktionslogik: Der Mann soll sich
\emph{qualifizieren}, um begehrt zu werden --- durch Status offline, durch
bezahlte Sichtbarkeit online. Beide Male ist die Qualifikation die
Eintrittskarte, nicht der Gewinn. Und beide Male ist der Ausstieg aus der
Qualifikation gleichbedeutend mit dem Ausscheiden aus dem Markt.

\section{Der Offline-Rest: Die unvermittelte Partnerwahl}

Bis hierher hat dieser Band fast ausschließlich den Online-Markt vermessen.
Das ist eine Verzerrung, die korrigiert werden muss: Ein erheblicher Teil der
Paare entsteht weiterhin ohne App --- und die Frage ist, ob dieser Rest anderen
Regeln gehorcht oder nur denselben Regeln in anderer Verpackung.

\subsection{Wo Paare heute entstehen}

Die beste Langzeitmessung stammt aus den USA: Michael Rosenfelds
HCMST-Erhebungen (\textit{How Couples Meet and Stay Together}), ausgewertet in
\textit{Disintermediating your friends} (PNAS 2019). \belegt\ Das Online-Kennenlernen
stieg von null (vor 1995) auf rund \textbf{39\,\%} aller neuen heterosexuellen
Paare (2017) und ist seither der mit Abstand häufigste einzelne Kennenlernweg.
Verdrängt wurde vor allem der klassische Vermittler: Das Kennenlernen
\emph{über Freunde} fiel von rund einem Drittel (Mitte der 1990er) auf etwa
ein Fünftel; auch Arbeitsplatz, Familie, Nachbarschaft und Kirche verloren
kontinuierlich. Für Deutschland zeigen Befragungsdaten (pairfam, FReDA,
Branchenerhebungen) dieselbe Richtung mit Verzögerung: In jüngeren Kohorten
lernt sich grob ein Drittel der neuen Paare online kennen, mit steigender
Tendenz --- und zugleich melden die Plattformen seit etwa 2024 Nutzerrückgänge
und eine dokumentierte \enquote{App-Müdigkeit} besonders der Jüngsten.
\plausibel

Der Offline-Rest ist also kein Relikt --- er ist die (knappe) Mehrheit der
Paarbildung, aber ein schrumpfender und strukturell geschwächter Kanal: Seine
klassischen Orte (Freundeskreis, Arbeitsplatz, Verein, Kneipe) verlieren an
Vermittlungskraft. Der Arbeitsplatz wird zusätzlich normativ stillgelegt, und
die \enquote{dritten Orte} im Sinne Oldenburgs --- Vereine, Stammlokale,
Gemeindestrukturen --- dünnen messbar aus. \plausibel\ Wer heute offline sucht,
sucht auf einem Feld, dessen Infrastruktur abgebaut wird.

\subsection{Die Regeln des Offline-Marktes: Nähe, Milieu, Zeit}

Die unvermittelte Partnerwahl gehorcht drei alten Gesetzen. \emph{Erstens der
Nähe} (Propinquity): Man verliebt sich in die, denen man wiederholt begegnet
--- der stärkste Einzelprädiktor der vormodernen Paarbildung. \emph{Zweitens
der Ähnlichkeit} (Homogamie): Paare gleichen sich in Bildung, Herkunft, Milieu
--- Bourdieus Habitus-Passung, die sich als \enquote{man versteht sich einfach}
anfühlt. \belegt\ \emph{Drittens der Zeit}: Offline-Bekanntschaft entsteht über
wiederholte Interaktion in geteilten Kontexten, und genau hier liegt der
entscheidende Unterschied zum App-Markt.

\begin{datenbox}[BEFUND · HUNT, EASTWICK \& FINKEL 2015, PSYCHOL.\ SCIENCE]
Paare, die sich erst kurz vor Beziehungsbeginn kennenlernten, waren einander
in physischer Attraktivität stark angeglichen (deutliche Matching-Korrelation).
Paare, die sich \emph{lange vor} Beziehungsbeginn kannten --- als Freunde,
Kollegen, Bekannte ---, zeigten dieses Matching kaum noch: Mit wachsender
Vertrautheit entkoppelt sich die Partnerwahl vom Marktwert des ersten
Eindrucks. Vertrautheit \enquote{ebnet das Spielfeld}.
\end{datenbox}

Das ist der theoretisch wichtigste Offline-Befund für diesen Band. Die
App-Selektion beruht fast vollständig auf dem ersten Eindruck --- Foto, Größe,
Beruf, also den Statusmarkern. Der Zeitkanal des Offline-Kennenlernens macht
erfahrbar, was kein Profil transportiert: Humor, Verlässlichkeit, Zuwendung,
Alltagstauglichkeit. Wo Vertrautheit wachsen darf, verliert die Statusordnung
einen Teil ihres Zugriffs.

\begin{gegen}
Daraus folgt nicht, dass offline die egalitäre Idylle wartet. Die Homogamie-Daten
zeigen das Gegenteil: Offline selektiert das \emph{Milieu}, bevor irgendjemand
selektieren kann. Wer im falschen Stadtteil wohnt, den falschen Beruf hat, im
falschen Verein ist, \emph{trifft} die Passenden nie --- die soziale Schließung
ersetzt den App-Filter und kann härter sein als er, weil sie unsichtbar bleibt.
Die Statuslogik verschwindet offline nicht; sie wandert von der expliziten
Auswahl in die implizite Vorauswahl der Begegnungsräume. Der Unterschied ist
real, aber er ist einer des Mechanismus, nicht der Existenz.
\end{gegen}

\begin{merksatz}
Der Offline-Weg hebt die Statuslogik nicht auf --- er verlangsamt sie. Und die
Verlangsamung ist genau der Kanal, den die App wegoptimiert hat: Zeit, in der
Vertrautheit den ersten Eindruck überschreiben kann.
\end{merksatz}

\subsection{Absichten, Commitment und die Älteren}

Bleibt die Absichtsfrage --- und mit ihr eine Hypothese, die sich vielen
männlichen Beobachtern aufdrängt: dass der Markt Optionsoffenhaltung belohne
und ernsthafte Bindungssuche bestrafe, und dass insbesondere Frauen sich
\enquote{alle Optionen offen} hielten. Diese Hypothese ist ernst zu nehmen ---
und ergebnisoffen zu prüfen.

Was Menschen \emph{angeben}: In der Pew-Erhebung 2023 nannten \textbf{44\,\%}
der aktiven Nutzer die Suche nach einer festen Partnerschaft als Hauptgrund,
40\,\% lockeres Daten, 24\,\% Gelegenheitssex --- Letzteres mit dem einzigen
signifikanten Geschlechterunterschied (Männer 31\,\%, Frauen 13\,\%); bei
allen anderen Motiven unterschieden sich die Geschlechter \emph{nicht}
signifikant. \belegt\ Nach Selbstauskunft ist der Beziehungswunsch also
mehrheitsfähig, beidgeschlechtlich --- und die Casual-Orientierung eher ein
männliches als ein weibliches Übergewicht. Die Hypothese in ihrer
geschlechtszugespitzten Form findet in den Absichtsdaten keine Stütze; unter
dem Vorbehalt von §\,3.4 gilt freilich: Das sind geäußerte, nicht gelebte
Absichten, verzerrt in beide Richtungen.

Was Menschen \emph{tun}: Hier trägt die Hypothese einen wahren Kern --- nur
ist er nicht geschlechtlich, sondern strukturell. Die Forschung zu Choice
Overload und zum \emph{Rejection Mindset} (Pronk \& Denissen 2020) zeigt: Je
mehr Optionen durchmustert werden, desto höher steigt die Ablehnungsquote und
desto geringer werden Zufriedenheit und Festlegungsbereitschaft --- ein
allgemein menschlicher Effekt, der bei Frauen im Online-Kontext steiler
ausfällt, weil das Überangebot ihn dort stärker füttert. \belegt\ Die
Optionsoffenhaltung ist also kein weiblicher Charakterzug, sondern eine
Funktion verfügbarer Optionen; die Marktmacht entscheidet, wer sie sich
leisten kann. Ghosting und \enquote{Situationships} sind die Alltagsformen
dieses Musters --- verbreitet, beidgeschlechtlich, und von den Betroffenen
beider Seiten als Bindungsscheu \emph{der anderen} erlebt. \plausibel

\begin{gegen}
Zur subjektiven Beobachtung gehört ihre Prüfung auf Wahrnehmungsverzerrung.
Die App-Ära macht Nicht-Bindung erstmals \emph{sichtbar}: Das offen gehaltene
Optionsportfolio war früher privat, heute ist es beobachtbare
Plattformmechanik. Sichtbarkeit ist nicht Zunahme. Zudem begegnet der ernsthaft
suchende Mann den bindungsbereiten Frauen seltener \emph{auf} den Apps ---
wer gefunden hat, verlässt den Markt; die Verbleibenden sind systematisch
selektiert (Survivorship). Beide Effekte lassen die Bindungsscheu größer
erscheinen, als die Verhaltensdaten sie ausweisen.
\end{gegen}

Und die Älteren: Das Segment 50plus ist das am schnellsten wachsende der
Partnersuche --- getrieben von Scheidung und Verwitwung ---, und es sucht nach
allen vorliegenden Befunden überwiegend Gefährtenschaft, nicht Abenteuer.
\plausibel\ Bemerkenswert ist die Marktmechanik: Im Alterssegment kippt das
Zahlenverhältnis. Durch die Lebenserwartungslücke und die männliche Präferenz
für jüngere Partnerinnen entsteht ein Frauenüberschuss --- die Marktmacht, die
in der Jugend bei den Frauen lag, wandert im Alter zu den Männern. Die
Statuslogik bleibt; nur ihre Nutznießer wechseln mit der Demografie. Für die
Kernthese ist das ein starker Beleg: Die Asymmetrie folgt der Knappheit, nicht
dem Geschlecht.

\section{Die Ausstiegsfrage}

Band VI schloss mit der Frage nach der Optionsvielfalt: Kann der Mann aus der
Funktion aussteigen? Die Antwort dort war ambivalent --- rechtlich ja,
ökonomisch, sozial und partnerschaftlich kaum. Der Partnermarkt schärft diese
Frage zu einer letzten: Ist der Ausstieg aus der \emph{Partnerschaft} der letzte
verbliebene Ausstieg aus der Funktion?

\subsection{Rückzug und Sexlosigkeit}

\begin{datenbox}[BEFUND · UEDA ET AL. 2020, JAMA NETWORK OPEN]
Der Anteil junger US-amerikanischer Männer (18--24) ohne sexuelle Aktivität im
zurückliegenden Jahr stieg über die 2000er und 2010er Jahre deutlich an --- in
manchen Erhebungen auf annähernd ein Drittel. Der Anstieg war bei Männern
stärker ausgeprägt als bei Frauen und korrelierte mit niedrigem Einkommen und
dem Wohnen bei den Eltern --- also mit \emph{Funktionsdefiziten}.
\end{datenbox}

Die Kopplung ist aufschlussreich: Nicht Sexlosigkeit erzeugt Funktionsverlust,
sondern Funktionsverlust (kein Einkommen, keine eigene Wohnung) erzeugt
Rückzug vom Partnermarkt. \plausibel\ Der Mann, der die Ernährerqualifikation
nicht erreicht, wird nicht abgelehnt --- er tritt gar nicht erst an. Das ist
Exklusion, die sich als Rückzug maskiert.

\subsection{Herbivoren, MGTOW und die Freiwilligkeitsfrage}

Kulturell haben sich Vokabeln für diesen Rückzug gebildet: die japanischen
\emph{s\=oshoku danshi} (\enquote{pflanzenfressende Männer}), die dem
Wettbewerb um Status, Karriere und Partnerin entsagen, und die westliche
MGTOW-Bewegung (\enquote{Men Going Their Own Way}), die den Ausstieg zur
Ideologie erhebt. \plausibel

Analytisch ist beides dieselbe Ambivalenz: Ist das souveräner Exit oder
rationalisierte Exklusion? Wer nicht gewählt wird und daraufhin erklärt, nicht
mehr wählen zu wollen, hat sich befreit --- oder seine Niederlage in eine
Entscheidung umgedeutet. Empirisch sind beide kaum zu trennen, weil dieselbe
Person beide Deutungen anbieten kann und die aufrichtigere oft selbst nicht
weiß, welche zutrifft.

\begin{merksatz}
Die Grenze zwischen \enquote{Ich will nicht} und \enquote{Ich werde nicht
gewollt} ist die unschärfste Linie dieses Bandes. Sie verläuft mitten durch die
Selbstwahrnehmung --- und wer sie zu klar zieht, lügt in die eine oder andere
Richtung.
\end{merksatz}

\subsection{Die Single-Gesellschaft und ihre technischen Nachkommen}

Der Rückzug ist längst statistischer Normalfall. \belegt\ Der Einpersonenhaushalt
ist in Deutschland die häufigste Haushaltsform; in Großstädten wie München oder
Berlin nähert sich sein Anteil der Hälfte aller Haushalte. Japan --- in vielem
ein Vorlaufszenario --- verzeichnet sinkende Heiratsquoten, eine
Rekord-Ehelosigkeit und eine Geburtenrate im Dauertief.

An dieser Stelle betritt die Provokation die Bühne, mit der dieser Band endet:
Wenn die Paarung selbst zur Minderheitenpraxis wird, welche Rolle übernehmen
dann die technischen Substitute? \spekulativ\ KI-Companions (Replika und
Nachfolger) bedienen bereits Millionen Nutzer mit parasozialer Zuwendung; die
\enquote{Einsamkeitsökonomie} wächst. Die zugespitzte Frage --- \enquote{Sind
Roboter unsere Nachkommen?} --- ist bewusst überzeichnet. Aber ihr ernster Kern
ist real: Wenn die Reproduktion an die Bedingung erfolgreicher Statusselektion
geknüpft bleibt und diese Selektion eine wachsende Männermehrheit ausschließt,
dann ist der Geburtenrückgang keine Laune, sondern eine \emph{Folge der
Funktionslogik} --- die Ordnung, die den Mann an die Funktion band, könnte an
eben dieser Bindung ihre eigene Fortsetzung verlieren.

\begin{gegen}
Vorsicht vor dem technikdeterministischen Kurzschluss. Sinkende Geburtenraten
haben viele Ursachen --- Bildungsexpansion, Verhütung, Urbanisierung, Kosten,
weibliche Autonomiegewinne, die ganz eigenständig und überwiegend als
Fortschritt zu werten sind. Die Selektionsasymmetrie des Partnermarkts als
\emph{die} Ursache auszurufen, überdehnt die Evidenz. Sie ist \emph{ein}
Faktor unter mehreren, und der KI-Companion ist bislang ein Nischenphänomen,
kein Reproduktionsersatz. Die Provokation taugt als Frage, nicht als Prognose.
\end{gegen}

\section{Synthese: Der Ausweg, der zurückführt}

Die Leitfrage dieses Bandes war, ob der Partnermarkt einen Ausweg aus der
Funktionslogik bietet oder ihr letzter Vollzugsraum ist. Die Antwort fällt
gegen den Ausweg aus.

Der Markt gibt dem Mann drei Optionen, und alle drei führen in die Funktion
zurück. \emph{Erstens} kann er sich qualifizieren --- Status erwerben, um
begehrt zu werden. Das ist die Unterwerfung unter die Funktion, nur unter
freieren Vorzeichen: kein Gesetz zwingt ihn, aber das Begehren tut es.
\emph{Zweitens} kann er in bezahlte Sichtbarkeit investieren --- der Zahlsklave,
der seinen Mangel abonniert. Das ist die Monetarisierung der Funktion.
\emph{Drittens} kann er sich zurückziehen --- Herbivore, MGTOW, freiwilliger
Single. Das ist der einzige echte Exit, aber er ist von der Exklusion kaum zu
unterscheiden und mündet in die Single-Gesellschaft, deren Fluchtpunkt die
technische Substitution ist.

\begin{merksatz}
Der Partnermarkt ist kein Gegenort zur Funktionslogik. Er ist der Ort, an dem
sie sich vom Zwang zum Begehren verflüssigt hat --- und dadurch unangreifbarer
wurde, weil niemand mehr einen Gesetzgeber anklagen kann. Man kann kein
Begehren vor Gericht bringen.
\end{merksatz}

Und die Begehrten? Sie sind, wie §\,4 zeigte, weder Täterinnen noch
Unbeteiligte. Sie vollstrecken eine Ordnung, die sie nicht entworfen haben, in
Akten, die sich als freier Geschmack erleben und es in gewissem Sinn auch sind.
Ihre Wahl belohnt Status, und diese Belohnung hält die Funktionslogik am Leben
--- aber sie tun es unter realen Anreizen (Kandiyoti), mit verinnerlichten
Schemata (Bourdieu), in einer Ordnung, die auch sie an ihre Kehrseite bindet.
Die Frage \enquote{Reproduzieren Frauen das Patriarchat?} ist mit einem
präzisen Doppel zu beantworten: \emph{Ja}, ihre aggregierte Wahl stabilisiert
die Statusordnung. Und \emph{nein}, das macht sie nicht schuldig, weil
Reproduktion einer Struktur keine Urheberschaft an ihr ist. Wer die Ordnung
vollstreckt, ohne sie gewählt zu haben, ist ihr Werkzeug, nicht ihr Autor.

Die redlichste Schlussfigur ist deshalb nicht Anklage, sondern eine geteilte
Unfreiheit: Auf diesem Markt begehren die einen nach oben und werden nicht
gehört, wählen die anderen von oben und tragen die Last, überhaupt gewählt
werden zu müssen, und über allen steht eine Ordnung, die sich in genau den
Gefühlen tarnt, die sich am freiesten anfühlen. Der Ausweg aus dem Patriarchat
führt nicht über den Partnermarkt. Er müsste über die Entkopplung von Status
und Partnerwert führen --- und ob die überhaupt denkbar ist, ohne die
zugrunde liegende ökonomische Asymmetrie zuerst aufzulösen, ist die Frage, an
der dieser Band bewusst nicht endet, sondern übergibt.

\subsection*{Ausblick: Band VIII}

Zwei Fortsetzungen bieten sich an, und der Cliffhanger dieses Bandes hält beide
offen:

\emph{Option A --- \enquote{Die Einsamkeitsökonomie}.} Wenn die Paarung zur
Minderheitenpraxis wird und die technischen Substitute wachsen: Was wird aus
Reproduktion, Bindung und Sorge in einer Gesellschaft, die das Begehren an eine
Selektion knüpft, die immer mehr Menschen ausschließt? Der KI-Companion als
Symptom, nicht als Lösung.

\emph{Option B --- \enquote{Der Staat als Ersatz-Patriarch}.} Wenn das Begehren
die Funktion nur noch unvollständig erzwingt, kehrt womöglich der alte
Erzwinger zurück: der Staat. Die Wiederkehr der Wehrpflicht-Debatte (das
Wehrdienst-Modernisierungsgesetz 2026, an das Band VI in §\,2 rührte) wäre dann
kein Zufall, sondern die Rückkehr der männlichen Pflicht durch die Vordertür,
nachdem sie durch die Hintertür des Begehrens nicht mehr zuverlässig
funktioniert.

Und eine dritte Tür, hier nur angelehnt: Was geschieht mit denen, die die
heterosexuelle Partnersuche ganz aufgeben --- und wohin gehen sie? Zwei
Richtungen wären zu prüfen, beide mit Vorsicht. Die \emph{Umorientierung}:
Sexuelle Orientierung ist nach vorherrschendem Forschungsstand keine
strategische Wahl; die dokumentierte Fluidität (vor allem bei Frauen, Diamond)
beschreibt gelebtes Erleben, kein Ausweichen vor einem schwierigen Markt ---
die Vorstellung einer \enquote{Konversion aus Marktfrust} hat keine
empirische Grundlage. \umstritten\ Und die \emph{Konvergenz in die
Incel-Radikalisierung}: die politisierte Form der Exklusion, in der die
Nicht-Wahl zur Identität und die Kränkung zur Ideologie wird. \spekulativ\
Band VIII wird diese Tür nur so weit öffnen, wie die Evidenz trägt.

Alle Fäden laufen auf dieselbe Frage zu, mit der die ganze Reihe begann und an
der sie weiterarbeitet: Was geschieht mit einer Ordnung, die Rechte an
Funktionen band, wenn die Funktionen sich auflösen --- aber die Bindung bleibt?

\section*{Quellen und Konfidenz}
\addcontentsline{toc}{section}{Quellen und Konfidenz}
\small
Die folgende Aufstellung nennt die tragenden Belege und ordnet ihre Belastbarkeit
ein. Vor Drucklegung sind die mit \umstritten\ und \spekulativ\ markierten
Stellen gegen die Primärquellen zu verifizieren; die App-Gini-Werte sind
ausdrücklich als Illustration, nicht als amtliche Statistik zu führen.

\begin{itemize}
\item \textbf{Rudder, C. (2014):} \textit{Dataclysm}. --- Bewertungsasymmetrie
(\textasciitilde80\,\%), Einkommenseffekt auf Anschreibhäufigkeit,
Alterspräferenzen. \belegt\ (als Plattformbefund) / \umstritten\ (als Aussage
über Präferenzen jenseits der Plattform).
\item \textbf{Tyson, G. et al. (2016):} \textit{A First Look at User Activity on
Tinder.} --- Match-Raten 0,6\,\% vs.\ 10\,\%. \belegt.
\item \textbf{Bruch, E. \& Newman, M. (2018):} \textit{Aspirational pursuit of
mates in online dating markets}, Science Advances. --- 25\,\%-Aspiration
(beidseitig), Desirability-Peaks nach Alter/Bildung. \belegt.
\item \textbf{Goldgeier, A. / Worst-Online-Dater:} Gini-Schätzungen (Hinge
\textasciitilde0,54 vs.\ 0,38; Tinder \textasciitilde0,58). \umstritten\
(Nicht-Peer-Review; als Größenordnung konsistent).
\item \textbf{Killewald, A. (2016):} \textit{Money, Work, and Marital
Stability}, ASR. --- Erwerbslosigkeit des Mannes erhöht Trennungsrisiko.
\belegt.
\item \textbf{Bertrand, M., Kamenica, E. \& Pan, J. (2015):} \textit{Gender
Identity and Relative Income within Households}, QJE. --- 50/50-Knick,
Kompensationsverhalten. \belegt.
\item \textbf{Ueda, P. et al. (2020):} \textit{Trends in Frequency of Sexual
Activity{\ldots}}, JAMA Network Open. --- Anstieg der Sexlosigkeit junger
Männer. \belegt\ (US-Daten; Übertragbarkeit auf D \umstritten).
\item \textbf{Bourdieu, P. (1998):} \textit{Die männliche Herrschaft.} ---
Symbolische Gewalt; der Herrschende als Beherrschter. \plausibel\ (Theorie).
\item \textbf{Kandiyoti, D. (1988):} \textit{Bargaining with Patriarchy.} ---
Patriarchaler Handel. \plausibel\ (Theorie).
\item \textbf{hooks, b. (2004):} \textit{The Will to Change.} --- Patriarchat
verstümmelt auch Männer. \plausibel\ (Theorie).
\item \textbf{Eagly, A. \& Wood, W. (1999)} vs.\ \textbf{Buss, D. \& Schmitt:}
Social Role Theory vs.\ Evolutionspsychologie; Gender-Equality-Paradox;
Eastwick \& Finkel (2008) zu stated vs.\ revealed preferences. \umstritten\
(offene Kontroverse).
\item \textbf{Esteve, A. et al.:} Umkehr der Bildungshypergamie. \belegt.
\item \textbf{Match Group / Marktstruktur:} Geschlechterverhältnis, Umsatz aus
zahlenden Männern, Super-Like-Erfolgsparadox. \plausibel\ (Branchendaten,
teils Schätzung).
\item \textbf{Destatis / BiB:} Einpersonenhaushalte als häufigste Haushaltsform;
Großstadtanteile. \belegt.
\item \textbf{KI-Companions / Einsamkeitsökonomie:} Replika u.\,a., parasoziale
Bindung. \spekulativ\ (in der Zuspitzung \enquote{Roboter als Nachkommen}).
\item \textbf{Toma, C., Hancock, J. \& Ellison, N. (2008):} \textit{Separating
Fact From Fiction}, PSPB. --- 81\,\% Abweichung bei Größe/Gewicht/Alter,
geschlechtstypische Richtung, Fotos am stärksten geschönt. \belegt.
\item \textbf{Rowatt, W. et al.:} Täuschungsbereitschaft steigt mit
Attraktivität des Gegenübers, ohne Geschlechtsunterschied. \plausibel.
\item \textbf{Pew Research Center (2023):} \textit{From Looking for Love to
Swiping the Field.} --- Motive der Nutzung: 44\,\% feste Partnerschaft,
24\,\% Gelegenheitssex (M 31\,\% / F 13\,\%, einziger signifikanter
Geschlechterunterschied). \belegt.
\item \textbf{Rosenfeld, M., Thomas, R. \& Hausen, S. (2019):}
\textit{Disintermediating your friends}, PNAS. --- Online \textasciitilde39\,\%
(2017), Kollaps der Freundeskreis-Vermittlung. \belegt\ (US-Daten; deutsche
Zahlen \plausibel).
\item \textbf{Hunt, L., Eastwick, P. \& Finkel, E. (2015):} \textit{Leveling
the Playing Field}, Psychol.\ Science. --- Vertrautheit entkoppelt Partnerwahl
vom Attraktivitäts-Matching. \belegt.
\item \textbf{Pronk, T. \& Denissen, J. (2020):} \textit{Rejection Mindset},
Psychol.\ Science. --- Ablehnungsquote steigt mit Optionsmenge,
geschlechtsübergreifend, online bei Frauen steiler. \belegt.
\item \textbf{Oldenburg, R.:} Dritte Orte; deutsches Vereins- und
Kneipensterben als Infrastrukturverlust der Offline-Paarbildung. \plausibel.
\item \textbf{Diamond, L.:} Sexuelle Fluidität (überwiegend bei Frauen
dokumentiert); keine Evidenz für \enquote{Konversion} als Marktausweichen.
\umstritten\ (für den Ausblick Band VIII).
\end{itemize}

\vspace{1em}
{\color{obsmuted}\rule{\textwidth}{0.4pt}}\\[0.4em]
{\ttfamily\scriptsize\color{obsmuted}
OBSERVATORY · Band VII \enquote{Der Markt der Begehrten} · Fortsetzung von
\enquote{Die Frau vor dem Recht} · Leitformel: Rechtsstellung folgt
Funktionsstellung · Diese Studie erhebt keinen Anspruch auf abgeschlossene
Wahrheit, sondern auf methodisch markierte Redlichkeit: Jede Konfidenzstufe ist
ausgewiesen, jede starke These trägt ihre Gegendarstellung.}
